\section{A fixed obstruction to source-only restrictions}
\label{sourceobs:section}

The mixed norm in Proposition~\ref{bisector:criterion} is an additional
geometric assumption. The following fixed example shows that it cannot
be supplied by the source dimensions, even after replacing the source
measure. The example lies strictly inside the already established
coherent dimension range. Its conclusion concerns the regularity of
the raw measure and is not a counterexample to positive distance length.

\begin{theorem}[Persistent source obstruction]\label{sourceobs:theorem}
There are a compact \(K\subset\R^2\) and a probability \(\mu\) on \(K\)
such that
\begin{equation}\label{sourceobs:dimensions}
 s=\HH K=\frac{11}{10},\qquad
 u=\PP K=\overline{\dim}_{\mathrm B}K=\frac{13}{11}<2s-1.
\end{equation}
The measure \(\mu\) is \(t\)-Frostman for every \(0<t<s\).
Let \(T(x)=x+(0,3)\), and let \(\nu\) be either \(\mu\) or
\(T_*\mu\). For every nonzero finite
positive measure \(\eta\) supported on \(K\), the joint distance law
\[
 \mathcal D_{\eta,\nu}(A)
 =\iint 1_A(y,|x-y|)\,d\eta(x)\,d\nu(y)
\]
has no density in \(L^{q,\infty}(\nu\times dt)\) for any \(q>9/8\).
In particular it has no \(L^2(\nu\times dt)\) density. For
\(\nu=\mu\), the conclusion already holds for the restriction
of the distance coordinate to \([1/64,1/16]\). Thus the full
pin measure can be supported on the same set \(K\), with the
obstruction occurring away from zero distance.
Separately, the full joint law with pin measure \(T_*\eta\)
has no weak-\(L^q\) density for any \(q>9/8\).

More generally, if \(\eta\le\mu\), \(\xi\le T_*\mu\), and \(\eta\) and
\((T^{-1})_*\xi\) have a common nonzero positive submeasure, the
corresponding joint law has no weak-\(L^q(\xi\times dt)\) density for
\(q>9/8\).

Consequently, if a nonzero positive \(\eta\) supported on \(K\) has
\(I_1(\eta)<\infty\), then for no \(1<\tau<s\) does its weighted
bisector measure have a density in
\(L^2_\theta L^{2/\tau}_c\). This applies to every nonzero
\(\eta\le\mu\), and to every Frostman source on \(K\) of exponent
greater than one.
\end{theorem}

The statement does not address arbitrary mutually singular source and
pin restrictions, or positive restrictions depending jointly on the
source and pin. In particular it does not preclude a joint density in
some \(L^p\) with \(p\) sufficiently close to one.

\subsection{A uniform finite block}

Set
\[
 \alpha=\frac2{11},\qquad \beta=\frac1{20},\qquad c=2^{-10}.
\]
For an integer \(L\ge10\), use the scales
\begin{equation}\label{sourceobs:scales}
 \delta=2^{-40L},\quad a=2^{-11L},\quad
 w=2^{-20L},\quad h=2^{-22L},\quad m=2^{-2L}.
\end{equation}
Thus \(w^2=\delta\), \(h=\delta^{s/2}\),
\(m=\delta^\beta\), and \(a=\delta^{\beta/\alpha}\).
Here \(a,m\) are geometric block parameters, independent of the
profile notation used in earlier sections.

Let \(B=2048\), and take the \(N_T=4^L=m^{-1}\) track centers
\[
 c_{(d_\ell)}
 =\frac18+\frac12\sum_{\ell=1}^L d_\ell B^{-\ell},
 \qquad d_\ell\in\{0,512,1024,1536\}.
\]
They have separation at least \(256a\). Uniformly in \(L\), their
normalized counting measure assigns mass at most \(CR^\alpha\) to
an interval of radius \(R\ge a\), and the centers can be covered by
at most \(CR^{-\alpha}\) intervals of length \(R\). To see this,
choose the base-\(B\) generation \(k\) with
\(B^{-k-1}<R\le B^{-k}\). An interval of this radius meets a
bounded number of generation-\(k\) cylinders, each of mass \(4^{-k}\);
there are \(4^k\) cylinders in total. The fixed contraction factor
\(1/2\) only changes the constants.

Put
\[
 N_S=2^{22L-4},\qquad N_X=2^{20L-4}.
\]
In track \(i\), for \(0\le k<N_S\) and \(0\le l<N_X\), place a
closed child square of side \(c\delta\) with center
\[
 z_{i,k,l}=
 \left(c_i+4\delta\left(l-\frac{N_X-1}{2}\right),
                  \frac14+kh\right).
\]
All squares are inside the unit square and have mutual separation
at least \(3\delta\). A track has horizontal extent at most \(w/4\).
The \(N_S\) slats in each track have vertical separation \(h\).
There are
\begin{equation}\label{sourceobs:branching}
 M_L=N_TN_SN_X=2^{44L-8}=2^{-8}\delta^{-s}
\end{equation}
children. Assign each child probability \(M_L^{-1}\). Then every
track has mass \(m\), and each of its slats has mass \(m/N_S\).
All the next estimates are independent of how the probability is
distributed inside each child.

\begin{lemma}[Uniform block bounds]\label{sourceobs:block}
For \(c\delta\le R\le1\), every such block measure satisfies
\[
 \mu_{\rm block}(B(x,R))\le CR^s,\qquad
 N(\supp\mu_{\rm block},R)\le CR^{-u},
\]
with constants independent of \(L\). At \(R=h\) one also has
\[
 \mu_{\rm block}(B(x,h))\le C\delta^{3s/2-1}.
\]
\end{lemma}
\begin{proof}
For \(\delta\le R\le1\), horizontal mass is at most \(CR^\alpha\)
for \(R\ge a\), \(Cm\) for \(w\le R\le a\), and \(CmR/w\)
for \(\delta\le R\le w\). Vertical mass is at most \(CR\) for
\(R\ge h\), and \(Ch\) for \(\delta\le R\le h\).
The product child structure and fixed child diameter therefore give
\[
\begin{array}{c|c|c}
\text{range}&\text{mass bound}&\text{cover bound}\\ \hline
a\le R\le1&CR^{1+\alpha}&CR^{-1-\alpha}\\
w\le R\le a&C\delta^\beta R&C\delta^{-\beta}R^{-1}\\
h\le R\le w&C\delta^{\beta-1/2}R^2
 &C\delta^{1/2-\beta}R^{-2}\\
\delta\le R\le h&C\delta^{s-1}R
 &C\delta^{1-s}R^{-1}.
\end{array}
\]
Every mass bound is at most \(CR^s\). Every cover bound is at
most \(CR^{-u}\), using
\[
 u=3-\frac2s,\qquad
 \frac{s}{2}(2-u)=1-\frac{s}{2}=\frac12-\beta.
\]
For \(c\delta\le R\le\delta\), only a bounded number of children
meet a ball, and a child needs a bounded number of radius-\(R\)
balls. The constants can depend on the fixed \(c\).
Finally, a radius-\(h\) ball meets a bounded number of tracks and
slats, and \(O(h/\delta)\) horizontal children, each of mass
comparable to \(\delta^s\). This gives the last estimate.
\end{proof}

\subsection{The recursive measure and exact dimensions}

Choose
\begin{equation}\label{sourceobs:depths}
 L_j=10+j^2\left(1+\sum_{k<j}L_k\right).
\end{equation}
Inside each stage-\((j-1)\) square insert the preceding \(L_j\)
block using its affine coordinates, with no rotation. Each child
is a similarity of ratio \(r_j=c\delta_j\). The intersection is
a compact set \(K\), and equal conditional child weights define a
probability \(\mu\). Write
\[
 \lambda_j=\prod_{k\le j}r_k,\qquad
 p_j=\prod_{k\le j}M_{L_k}^{-1},\qquad P_j=p_j^{-1},
 \qquad \lambda_0=p_0=1.
\]
There are \(P_j\) stage-\(j\) squares, each of side \(\lambda_j\)
and mass \(p_j\). Their separation is at least a fixed large
multiple of \(\lambda_j\): this follows for siblings from
\(3\delta_j/(c\delta_j)=3/c\), and for different parents by
induction. Thus a ball of radius at most \(\lambda_j\) meets
only a bounded number of stage-\(j\) squares. In particular every
point of \(K\) has an unambiguous parent, track, slat, and child
at every stage.

For \(\lambda_j\le r\le\lambda_{j-1}\), applying
Lemma~\ref{sourceobs:block} in each parent met by the ball gives
\begin{equation}\label{sourceobs:frostman}
 \mu(B(x,r))\le
 C p_{j-1}\lambda_{j-1}^{-s}r^s.
\end{equation}
The factor \(p_{j-1}\lambda_{j-1}^{-s}\) equals
\((2^8c^{-s})^{j-1}\). For every \(0<t<s\),
\[
 p_{j-1}\lambda_{j-1}^{-s}r^{s-t}
 \le (2^8c^{-s})^{j-1}\lambda_{j-1}^{s-t},
\]
which is uniformly bounded. Hence \(\mu\) is \(t\)-Frostman for
every \(0<t<s\), and \(\HH K\ge s\). Conversely the level covers
have
\[
 \frac{\log P_j}{-\log\lambda_j}\longrightarrow s,
\]
so \(\HH K\le s\).

For the covering bound, the same lemma yields
\[
 N(K,r)\le C P_{j-1}(r/\lambda_{j-1})^{-u}.
\]
The prefactor is bounded, since
\[
 P_{j-1}\lambda_{j-1}^u
 =\prod_{k<j}(2^{-8}c^u\delta_k^{u-s})\le1.
\]
Thus \(\overline{\dim}_{\mathrm B}K\le u\) and \(\PP K\le u\).

For the reverse inequality set \(r_j^*=\lambda_{j-1}h_j\).
The refined block estimate gives
\[
 \sup_x\mu(B(x,r_j^*))
 \le Cp_{j-1}\delta_j^{3s/2-1}.
\]
Consequently every \(A\subset K\) with \(\mu^*(A)>0\) satisfies
\[
 N(A,r_j^*)\ge
 \frac{\mu^*(A)}{Cp_{j-1}\delta_j^{3s/2-1}}.
\]
Since \(L_j/\sum_{k<j}L_k\to\infty\), logarithms imply
\[
 \overline{\dim}_{\mathrm B}A
 \ge\frac{3s/2-1}{s/2}=u.
\]
Every countable cover of \(K\) contains a piece of positive
\(\mu\) outer measure. The countable-cover characterization of
packing dimension therefore gives \(\PP K\ge u\), completing
\eqref{sourceobs:dimensions}.

\subsection{Collisions that persist under source restriction}

We first use the translated pin measure \(\nu=T_*\mu\).
Fix any finite positive \(\eta\) supported on \(K\), of mass
\(M>0\). At stage \(j\), abbreviate
\[
 \lambda=\lambda_{j-1},\quad p=p_{j-1},\quad
 \delta=\delta_j,\quad h=h_j,\quad w=w_j,\quad m=m_j.
\]
For each parent \(Q\), track \(i\), and slat \(k\), put
\[
 q_{Q,i,k}=\eta(\text{slat }(Q,i,k)),\qquad
 q_{Q,i}=\sum_kq_{Q,i,k},\qquad M_Q=\sum_iq_{Q,i}.
\]
Here a slat means the union of all its descendants in \(K\).
These are disjoint compact sets.

For a pin \(y\) in the translated track \(T(Q,i)\), any two
source points \(x,x'\) in slat \((Q,i,k)\) obey
\begin{equation}\label{sourceobs:distancewidth}
 \big||x-y|-|x'-y|\big|
 \le C(\lambda\delta+\lambda^2w^2)
 \le C\lambda\delta.
\end{equation}
The source vertical coordinates differ by \(O(\lambda\delta)\),
and their horizontal displacements from \(y\) are
\(O(\lambda w)\). Apply the difference-of-squares identity,
using the fixed positive separation of \(K\) and \(T(K)\).
Corresponding parent copies are separated by the global translation
3, not by \(3\lambda\); this explains the favorable
\(\lambda^2w^2\) term.

Choose nonnegative \(\varphi\in C_c^\infty(\R)\) of integral one.
Its autocorrelation is positive near zero. For
\(\varepsilon_j=C_0\lambda\delta\), with \(C_0\) a sufficiently
large fixed constant, \eqref{sourceobs:distancewidth} implies
\[
 \int|\varphi_{\varepsilon_j}*(d_y)_*\eta(t)|^2\,dt
 \ge\frac{c_1}{\lambda\delta}\sum_kq_{Q,i,k}^2
 \quad(y\in T(Q,i)).
\]
The pin mass of this parent-track is \(pm\). Since
\(N_S^{-1}=16h\), successive applications of Cauchy--Schwarz
over slats, tracks, and parents give
\begin{align}
 \int\|\varphi_{\varepsilon_j}*(d_y)_*\eta\|_2^2\,d\nu(y)
 &\ge\frac{c_1pmh}{\lambda\delta}\sum_{Q,i}q_{Q,i}^2
 \notag\\
 &\ge\frac{c_1pm^2h}{\lambda\delta}\sum_QM_Q^2
 \notag\\
 &\ge c_1M^2p^2\lambda^{-1}\delta^{3s/2-2}.
 \label{sourceobs:collision}
\end{align}
At \(s=11/10\), the last exponent is \(-7/20\), and
\[
 \log_2(p_{j-1}^2\lambda_{j-1}^{-1}\delta_j^{-7/20})
 =14L_j-48\sum_{k<j}L_k+26(j-1)\longrightarrow+\infty.
\]
Thus no joint \(L^2\) density exists, since convolution would
contract its \(L^2\) norm.

For pins \(T_*\eta\), the pin mass in the corresponding parent-track
is \(q_{Q,i}\). Jensen's inequality over the \(1/(pm)\)
parent-tracks gives instead
\[
 \frac{c_1h}{\lambda\delta}\sum_{Q,i}q_{Q,i}^3
 \ge c_1M^3p^2\lambda^{-1}\delta^{3s/2-2}\longrightarrow\infty.
\]
These estimates require only that \(\eta\) be supported on \(K\);
no domination by the natural measure is used.

\subsection{Failure of weak \(L^q\) above \(9/8\)}

For each \(y\in T(Q,i)\), take \(S_j(y)\) to be the union over
slats \(k\) of intervals of length \(C\lambda\delta\) centered
at the distances from \(y\) to the corresponding source slat
centers. The constant is fixed and large enough that
\eqref{sourceobs:distancewidth}, also with a slat center in place
of one source point, puts all distances from the source
parent-track inside this union. The finitely many centers and
compact pieces show that \(S_j=\{(y,t):t\in S_j(y)\}\) is Borel.
For \(\omega=\nu\times dt\),
\begin{equation}\label{sourceobs:support}
 \omega(S_j)\le C\lambda\delta^{1-s/2},\qquad
 \mathcal D_{\eta,\nu}(S_j)\ge
 \sum_{Q,i}q_{Q,i}\nu(T(Q,i))=pmM.
\end{equation}

For \(q>1\), the distribution-function formula gives
\[
 \int_A F\,d\omega
 \le\frac{q}{q-1}\|F\|_{L^{q,\infty}(\omega)}
                  \omega(A)^{1-1/q}
 \qquad(F\ge0).
\]
Thus a weak-\(L^q\) density of the joint law would satisfy
\begin{equation}\label{sourceobs:weaklower}
 \|F\|_{L^{q,\infty}}\ge
 c_qMp\lambda^{-(1-1/q)}
 \delta^{\beta-(1-s/2)(1-1/q)}.
\end{equation}
The exponent of \(\delta\) is negative precisely for
\[
 q>\frac{2-s}{3-2s}=\frac98.
\]
For each fixed such \(q\), the stage choice
\eqref{sourceobs:depths} makes the right side of
\eqref{sourceobs:weaklower} diverge, because the contribution of
\(L_j\) dominates every earlier depth. This proves the asserted
weak-\(L^q\) failure. We make no claim at \(q=9/8\).

For pins \(T_*\eta\), the product measure of \(S_j\) is at most
\(CM\lambda\delta^{1-s/2}\), whereas its joint distance mass is
at least
\[
 \sum_{Q,i}q_{Q,i}^2\ge M^2pm.
\]
The same exponent therefore proves the matched-pin statement.
If a nonzero \(\zeta\le\eta\) also satisfies \(T_*\zeta\le\xi\),
the actual source--pin law dominates the law for source \(\zeta\)
and pins \(T_*\zeta\). Its mass on \(S_j\) is at least
\(\zeta(K)^2pm\), while
\[
 (\xi\times dt)(S_j)
 \le C\xi(\R^2)\lambda\delta^{1-s/2}.
\]
Applying the same support inequality relative to \(\xi\times dt\)
proves the common-submeasure variant. No change of base measure
in a Lorentz norm is being asserted here.

\subsection{The same-set full pin measure}

We now use \(\nu=\mu\), and keep the relevant distances in a
fixed interval separated from zero. The first-stage number \(N_S\)
is even. Define \(F:K\to K\) by replacing the first vertical slat
digit \(k\) with \(k+N_S/2\) modulo \(N_S\), leaving the first
track and horizontal-child indices, and every subsequent digit,
unchanged. Unique coding makes \(F\) a Borel involution, and equal
conditional weights give \(F_*\mu=\mu\).

On each first-stage child, \(F\) is a translation by either
\((0,1/32)\) or \((0,-1/32)\), since \(N_Sh_1/2=1/32\).
For every \(j\ge2\), it maps each stage-\((j-1)\) parent \(Q\)
onto a parent \(F(Q)\) by that same translation, preserving
all relative track, slat, and child geometry. In particular
\(\mu(F(Q,i))=pm\).

For \(y\in F(Q,i)\) and \(x,x'\) in source slat \((Q,i,k)\),
the proof of \eqref{sourceobs:distancewidth} applies with
fixed separation \(1/32\) instead of \(3\). Indeed
\(\lambda\le\lambda_1\) is far smaller than \(1/64\), and
\[
 \left||x-y|-\frac1{32}\right|\le\sqrt2\,\lambda.
\]
Thus their distance oscillation is at most \(C\lambda\delta\),
and these distances and their \(C\lambda\delta\) neighborhoods
lie in \(J=[1/64,1/16]\).

For \(y\in F(Q,i)\), form \(S_j(y)\) using the distances to
all source slat centers in \(Q,i\), exactly as in the preceding
subsection. The parent correspondence is a bijection, so
\begin{equation}\label{sourceobs:samepins}
 \begin{gathered}
 S_j\subset K\times J,\qquad
 (\mu\times dt)(S_j)\le C\lambda\delta^{1-s/2},\\
 \mathcal D_{\eta,\mu}(S_j)\ge pmM.
 \end{gathered}
\end{equation}
The support inequality giving \eqref{sourceobs:weaklower} now
applies even if only the restriction of the law to \(K\times J\)
were assumed to have a weak-\(L^q\) density. This proves the
same-set full-pin assertion away from the radial diagonal.
The collision proof \eqref{sourceobs:collision} likewise applies
with \(F(Q,i)\) in place of \(T(Q,i)\).

This argument does not assert that an arbitrary independently
selected pin restriction retains the requisite matching
parent-track masses. The translated matched-pin conclusion was
proved separately above.

\subsection{Consequence for the bisector criterion}

Let \(1<\tau<s\), and let \(\eta\) be a nonzero positive source
on \(K\) with finite first energy. Choosing \(\tau<t<s\), the
\(t\)-Frostman estimate gives \(I_\tau(\mu)<\infty\), so we may
use the full same-set pin measure \(\nu=\mu\).
If \(\mathcal B_\eta\) had density in
\(L^2_\theta L^{2/\tau}_c\), Proposition~\ref{bisector:criterion}
would give a raw \(L^2(\mu\times dt)\) joint distance density,
contrary to \eqref{sourceobs:samepins}. Every submeasure of \(\mu\),
and every Frostman source on \(K\) of exponent greater than one,
has the required finite first energy. This completes the proof of
Theorem~\ref{sourceobs:theorem}.

The strict inequality \(13/11<6/5=2s-1\) is deliberate. It shows
that this mixed norm is not even forced throughout the established
coherent range. The obstruction leaves open positive distance length,
absolute continuity or smaller \(L^p\) exponents, and selections
depending jointly on the source and pin.
